%======================================================================
\section{Relative entropy}\label{sec:relent}
%======================================================================

\begin{objectives}
  You will compute the classical and quantum relative entropies, prove Gibbs' and Klein's inequalities, and use them to prove the main properties of entropies: the non-negativity of the mutual information, subadditivity, concavity, and the fact that measuring in a basis increases entropy.
\end{objectives}

\begin{recall}
  The von Neumann entropy (Section~\ref{sec:vonneumann}) and Shannon entropy, Jensen's inequality (Section~\ref{sec:shannon}); the partial trace and $\Tr((M\otimes\Id)\rho^{AB})=\Tr(M\rho^A)$ (Section~\ref{sec:density}); projectors and the functional calculus (Section~\ref{sec:operators}).
\end{recall}

How different are two probability distributions? Two quantum states? The entropy of a single distribution cannot answer this, but there is a natural asymmetric measure of how surprised we are to see data generated by $p$ when we believed it was generated by $q$: the relative entropy. It is also the tool that gives the cleanest proofs of nearly every inequality in this chapter.

\subsection{The classical relative entropy}

\begin{definition}[Relative entropy of distributions]\label{def:kl}
  For probability distributions $p,q$ on the same finite set, the \term{relative entropy} (or Kullback--Leibler divergence) is
  \[
    \KL pq=\sum_xp(x)\log\frac{p(x)}{q(x)},
  \]
  with the conventions $0\log\frac0q=0$ and $p\log\frac p0=+\infty$ for $p>0$.
\end{definition}

It is the average, under $p$, of the \emph{extra} surprise $\log\frac{p}{q}$ of $q$ over $p$. It is not symmetric, and it is not a distance, but it is a faithful measure of discrepancy, by the following inequality.

\begin{theorem}[Gibbs' inequality]\label{thm:gibbs}
  For all distributions $p,q$, $\ \KL pq\ge0$, and $\KL pq=0$ iff $p=q$. More generally the inequality holds if $q$ is only a \emph{sub-normalised} distribution ($q(x)\ge0$, $\sum_xq(x)\le1$).
\end{theorem}
\begin{proof}
  If $p(x)>0$ for some $x$ with $q(x)=0$, then $\KL pq=\infty$ and there is nothing to prove, so assume $\supp p\subseteq\supp q$.
  \begin{steps}
    \item Use the inequality $\ln t\le t-1$, with equality iff $t=1$.
    \item \emph{Apply it to $t=q(x)/p(x)$ for each $x$ with $p(x)>0$:} $-\KL pq=\frac1{\ln2}\sum_{p(x)>0}p(x)\ln\frac{q(x)}{p(x)}\le\frac1{\ln2}\sum_{p(x)>0}p(x)\Bigl(\frac{q(x)}{p(x)}-1\Bigr)$.
    \item \emph{Simplify:} the right-hand side is $\frac1{\ln2}\bigl(\sum_{p(x)>0}q(x)-1\bigr)\le0$, since $\sum q\le1$.
    \item \emph{Equality:} we need $q(x)=p(x)$ whenever $p(x)>0$ and $\sum_{p(x)>0}q(x)=1$, which together force $q=p$. \qedhere
  \end{steps}
\end{proof}

Gibbs' inequality gives two earlier results at once, and a third one that we had postponed.

\begin{corollary}\label{cor:gibbs-consequences}
  \begin{enumerate}
    \item[(a)] $\Hs(p)\le\log n$, with equality iff $p$ is uniform.
    \item[(b)] $I(X\!:\!Y)=\KL{p_{XY}}{p_Xp_Y}\ge0$, with equality iff $X,Y$ are independent.
    \item[(c)] Conditioning reduces entropy: $\Hs(Y\mid X)\le\Hs(Y)$.
  \end{enumerate}
\end{corollary}
\begin{proof}
  \begin{steps}
    \item (a) Take $q$ uniform: $\KL pq=\sum_xp(x)\log(np(x))=\log n-\Hs(p)\ge0$.
    \item (b) Expanding, $\KL{p_{XY}}{p_Xp_Y}=\sum p(x,y)\bigl[\log p(x,y)-\log p_X(x)-\log p_Y(y)\bigr]=-\Hs(X,Y)+\Hs(X)+\Hs(Y)=I(X\!:\!Y)$. Gibbs gives $\ge0$, with equality iff $p(x,y)=p_X(x)p_Y(y)$.
    \item (c) $\Hs(Y)-\Hs(Y\mid X)=I(X\!:\!Y)\ge0$. \qedhere
  \end{steps}
\end{proof}

\begin{example}\label{ex:kl}
  Compute $\KL pq$ and $\KL qp$ for $p=(\tfrac12,\tfrac12)$ and $q=(\tfrac14,\tfrac34)$.
\end{example}
\begin{solution}
  \begin{steps}
    \item $\KL pq=\tfrac12\log\frac{1/2}{1/4}+\tfrac12\log\frac{1/2}{3/4}=\tfrac12(1)+\tfrac12(-0.585)\approx0.208$ bits.
    \item $\KL qp=\tfrac14\log\frac{1/4}{1/2}+\tfrac34\log\frac{3/4}{1/2}=\tfrac14(-1)+\tfrac34(0.585)\approx0.189$ bits.
  \end{steps}
  Both are positive, as Gibbs says, but they differ: relative entropy is not symmetric.
\end{solution}

\begin{lookahead}
  Relative entropy is the exact measure of how well a test can tell $p$ from $q$. By Stein's lemma, given $n$ independent samples, any test that wrongly accepts $q$ with probability bounded away from $1$ when the truth is $p$ can make its other error shrink at best like $2^{-n\,\KL pq}$. The quantum version, below, plays the same role for distinguishing quantum states, and extends the work of Section~\ref{sec:measurement} on distinguishing states with a single measurement.
\end{lookahead}

\subsection{The quantum relative entropy}

\begin{definition}[Quantum relative entropy]\label{def:qre}
  For density operators $\rho,\sigma$, the \term{relative entropy} of $\rho$ with respect to $\sigma$ is
  \[
    \Srel\rho\sigma=\Tr\bigl(\rho\log\rho\bigr)-\Tr\bigl(\rho\log\sigma\bigr)\quad\text{if }\supp\rho\subseteq\supp\sigma,\qquad\text{and }+\infty\text{ otherwise}.
  \]
  Here $\log\sigma$ is the operator logarithm on the support of $\sigma$ (Definition~\ref{def:functional}).
\end{definition}

If $\rho$ and $\sigma$ commute, they share an eigenbasis and $\Srel\rho\sigma$ is just the classical $\KL pq$ for the two lists of eigenvalues. The quantum case adds the possibility that they do \emph{not} commute.

\begin{example}\label{ex:qre}
  Compute $\Srel\rho\sigma$ for: (a) $\rho=\ket0\!\bra0$, $\sigma=\Id/2$; (b) $\rho=\ket0\!\bra0$, $\sigma=\ket1\!\bra1$; (c) $\rho=\diag(\tfrac34,\tfrac14)$, $\sigma=\Id/2$; (d) $\rho=\diag(\tfrac34,\tfrac14)$ and $\sigma=\tfrac34\ket+\!\bra++\tfrac14\ket-\!\bra-$.
\end{example}
\begin{solution}
  \begin{steps}
    \item (a) $\Tr\rho\log\rho=0$ and $\log\sigma=-\Id$ (in bits), so $\Srel\rho\sigma=0-(-1)=1$ bit.
    \item (b) $\supp\rho=\Span\{\ket0\}\not\subseteq\supp\sigma=\Span\{\ket1\}$, so $\Srel\rho\sigma=+\infty$: a state that is certainly $\ket0$ is infinitely surprising if we believed $\ket1$.
    \item (c) $\Srel\rho\sigma=-\Ss(\rho)-\Tr\rho\log\sigma=-0.811+1=0.189$ bits, which is $\log2-\Ss(\rho)$.
    \item (d) The operators do not commute. We use $\log\sigma=\log\tfrac34\,\ket+\!\bra++\log\tfrac14\,\ket-\!\bra-$, so $\Tr\rho\log\sigma=\log\tfrac34\bra+\rho\ket++\log\tfrac14\bra-\rho\ket-$. Since $\bra\pm\rho\ket\pm=\tfrac12(\tfrac34+\tfrac14)=\tfrac12$, this is $\tfrac12\log\tfrac34+\tfrac12\log\tfrac14\approx-0.208-1=-1.208$. Hence $\Srel\rho\sigma=-0.811+1.208=0.396$ bits.
  \end{steps}
\end{solution}

The fundamental result is the quantum Gibbs inequality, due to Klein.

\begin{theorem}[Klein's inequality]\label{thm:klein}
  $\Srel\rho\sigma\ge0$, with equality iff $\rho=\sigma$.
\end{theorem}
\begin{proof}
  Assume $\supp\rho\subseteq\supp\sigma$ (otherwise the left side is $+\infty$). Write $\rho=\sum_ip_i\ket i\!\bra i$ and $\sigma=\sum_jq_j\ket{j'}\!\bra{j'}$ with two orthonormal bases, in general different. Let $P_{ij}=\abs{\braket i{j'}}^2$.
  \begin{steps}
    \item \emph{The matrix $P$.} $\sum_jP_{ij}=1$ for each $i$ (completeness of $\{\ket{j'}\}$), and $\sum_iP_{ij}=1$ for each $j$. So the rows and columns of $P$ are probability distributions: $P$ is \emph{doubly stochastic}.
    \item \emph{Rewrite the relative entropy.} $\Tr\rho\log\sigma=\sum_ip_i\bra i\log\sigma\ket i=\sum_ip_i\sum_jP_{ij}\log q_j$ (sums over $q_j>0$; the terms with $p_i>0$ only involve this part of $\sigma$ because $\supp\rho\subseteq\supp\sigma$). Hence
      \[
        \Srel\rho\sigma=\sum_ip_i\log p_i-\sum_ip_i\sum_jP_{ij}\log q_j .
      \]
    \item \emph{Jensen.} For fixed $i$, $(P_{ij})_j$ is a distribution and $\log$ is concave, so $\sum_jP_{ij}\log q_j\le\log\sum_jP_{ij}q_j=\log r_i$, where $r_i=\sum_jP_{ij}q_j$.
    \item \emph{A classical divergence appears.} Therefore $\Srel\rho\sigma\ge\sum_ip_i\log\frac{p_i}{r_i}=\KL pr$.
    \item \emph{$r$ is sub-normalised.} $\sum_ir_i=\sum_jq_j\sum_iP_{ij}=\sum_jq_j=1$ (the restriction to $q_j>0$ only makes it smaller, never larger than $1$).
    \item \emph{Conclude.} By Gibbs' inequality (Theorem~\ref{thm:gibbs}), $\KL pr\ge0$, so $\Srel\rho\sigma\ge0$.
    \item \emph{Equality.} If $\Srel\rho\sigma=0$, then equality holds in Jensen: for each $i$ with $p_i>0$, $q_j$ is constant on those $j$ with $P_{ij}>0$, equal to $r_i$; and $p=r$. Hence $\sigma\ket i=\sum_jq_j\braket{j'}{i}\ket{j'}=p_i\ket i$ for each $i$ with $p_i>0$. These eigenvectors already account for trace $\sum_ip_i=1$ of $\sigma$, so $\sigma$ has no further non-zero eigenvalue and $\sigma=\rho$. \qedhere
  \end{steps}
\end{proof}

\subsection{Consequences}

Klein's inequality quickly produces the properties promised in Section~\ref{sec:vonneumann}.

\begin{corollary}[Properties of the von Neumann entropy]\label{cor:vn-consequences}
  \begin{enumerate}
    \item[(a)] $\Ss(\rho)\le\log d$, with equality iff $\rho=\Id/d$.
    \item[(b)] \emph{Subadditivity:} $\Ss(\rho^{AB})\le\Ss(\rho^A)+\Ss(\rho^B)$, with equality iff $\rho^{AB}=\rho^A\otimes\rho^B$. Equivalently, the quantum mutual information $I(A\!:\!B)=\Ss(\rho^A)+\Ss(\rho^B)-\Ss(\rho^{AB})=\Srel{\rho^{AB}}{\rho^A\otimes\rho^B}\ge0$.
    \item[(c)] \emph{Concavity:} $\Ss\bigl(\sum_ip_i\rho_i\bigr)\ge\sum_ip_i\Ss(\rho_i)$.
    \item[(d)] \emph{Measuring in a basis increases entropy:} if $\mathcal D(\rho)=\sum_i\ket i\!\bra i\rho\ket i\!\bra i$ is the state after a non-selective measurement in an orthonormal basis $\{\ket i\}$, then $\Ss(\mathcal D(\rho))=\Ss(\rho)+\Srel\rho{\mathcal D(\rho)}\ge\Ss(\rho)$.
  \end{enumerate}
\end{corollary}
\begin{proof}
  \begin{steps}
    \item (a) Take $\sigma=\Id/d$: $\Srel\rho{\Id/d}=\Tr\rho\log\rho-\Tr\rho\log\frac{\Id}d=-\Ss(\rho)+\log d\ge0$, with equality iff $\rho=\Id/d$.
    \item (b) \emph{The logarithm of a product.} If $\rho^A=\sum_ap_a\ket a\!\bra a$ and $\rho^B=\sum_bq_b\ket b\!\bra b$ then $\rho^A\otimes\rho^B$ has eigenvectors $\ket a\ket b$ with eigenvalues $p_aq_b$, so $\log(\rho^A\otimes\rho^B)=\log\rho^A\otimes\Id+\Id\otimes\log\rho^B$. \emph{Take the expectation} in $\rho^{AB}$ using Proposition~\ref{prop:reduced}: $\Tr\bigl(\rho^{AB}(\log\rho^A\otimes\Id)\bigr)=\Tr\rho^A\log\rho^A=-\Ss(\rho^A)$, and likewise for $B$. Hence $\Srel{\rho^{AB}}{\rho^A\otimes\rho^B}=-\Ss(\rho^{AB})+\Ss(\rho^A)+\Ss(\rho^B)$, which is $\ge0$ by Klein, with equality iff $\rho^{AB}=\rho^A\otimes\rho^B$.
    \item (c) Let $\bar\rho=\sum_ip_i\rho_i$. Since $\supp\rho_i\subseteq\supp\bar\rho$, $\sum_ip_i\Srel{\rho_i}{\bar\rho}=\sum_ip_i\bigl(-\Ss(\rho_i)\bigr)-\Tr\bar\rho\log\bar\rho=\Ss(\bar\rho)-\sum_ip_i\Ss(\rho_i)$, and each term of the left side is $\ge0$.
    \item (d) $\log\mathcal D(\rho)$ is diagonal in $\{\ket i\}$, and for any diagonal $D$ one has $\Tr(\rho D)=\Tr(\mathcal D(\rho)D)$. Hence $\Tr\rho\log\mathcal D(\rho)=\Tr\mathcal D(\rho)\log\mathcal D(\rho)=-\Ss(\mathcal D(\rho))$. So $\Srel\rho{\mathcal D(\rho)}=-\Ss(\rho)+\Ss(\mathcal D(\rho))\ge0$. \qedhere
  \end{steps}
\end{proof}

\begin{example}\label{ex:entropy-increase}
  Verify (d) for $\rho=\ket+\!\bra+$ with the computational basis.
\end{example}
\begin{solution}
  \begin{steps}
    \item \emph{Before:} $\rho$ is pure, so $\Ss(\rho)=0$.
    \item \emph{After:} $\mathcal D(\rho)=\tfrac12\ket0\!\bra0+\tfrac12\ket1\!\bra1=\Id/2$, with $\Ss=1$.
    \item \emph{The relative entropy:} $\Srel\rho{\mathcal D(\rho)}=1-0=1$, consistent with the identity of (d).
  \end{steps}
  A projective measurement that we forget the result of destroys the coherence between $\ket0$ and $\ket1$ and raises the entropy from $0$ to its maximum.
\end{solution}

\begin{example}[Distinguishability can only be lost by discarding]\label{ex:monotone}
  Compare $\Srel{\rho^{AB}}{\sigma^{AB}}$ for $\rho^{AB}=\ket{\Phi^+}\!\bra{\Phi^+}$, $\sigma^{AB}=\Id/4$ with the relative entropy of the reduced states.
\end{example}
\begin{solution}
  \begin{steps}
    \item \emph{Joint:} $\Srel{\rho^{AB}}{\Id/4}=\log4-\Ss(\rho^{AB})=2-0=2$ bits.
    \item \emph{Reduced:} $\rho^A=\Id/2=\sigma^A$, so $\Srel{\rho^A}{\sigma^A}=0$.
  \end{steps}
  Two bits of distinguishability between the global states vanish when we look at one qubit only: all of it lived in the correlations. In general $\Srel{\rho^A}{\sigma^A}\le\Srel{\rho^{AB}}{\sigma^{AB}}$ (\emph{monotonicity under partial trace}), a deep theorem of Lieb, Ruskai and Lindblad, which gives the strong subadditivity of the entropy and says that no physical operation can increase distinguishability. It is proved in~\cite{wilde2017}.
\end{solution}

\begin{lookahead}
  Relative entropy is the unifying quantity. The mutual information is a relative entropy to a product; the entanglement of a state can be quantified by its relative entropy to the closest unentangled state; and in thermodynamics, with $\sigma=e^{-\beta H}/Z$ the Gibbs state, $\Srel\rho\sigma=\beta\bigl(F(\rho)-F(\sigma)\bigr)$ is the excess free energy of $\rho$ (Section~\ref{sec:blackbody} uses this identity to show that the Gibbs state has the largest entropy at fixed energy). It is also the exponent that governs the optimal discrimination of two quantum states from many copies (the quantum Stein lemma), and the natural next step after the one-shot discrimination of Section~\ref{sec:measurement}. See Vedral's review~\cite{vedral2002}.
\end{lookahead}

\begin{exerciseset}
  \exercise{1} Compute $\KL pq$ and $\KL qp$ for $p=(0.9,0.1)$, $q=(0.5,0.5)$, and for $p=(\tfrac12,\tfrac12,0)$, $q=(\tfrac13,\tfrac13,\tfrac13)$ and the reverse.

  \exercise{2} Prove $\ln t\le t-1$ for $t>0$ (study $g(t)=t-1-\ln t$), with equality only at $t=1$, and use it to give a one-line proof that $\KL pq\ge0$ for sub-normalised $q$.

  \exercise{2} Show that the mutual information can be written $I(X\!:\!Y)=\sum_{x,y}p(x,y)\log\frac{p(x,y)}{p_X(x)p_Y(y)}$, and compute it for the joint distribution of Exercise~\ref{ex:joint-dist}.

  \exercise{2} For the binary distributions $p=(a,1-a)$ and $q=(b,1-b)$ write $\KL pq$ explicitly, and show it vanishes only for $a=b$.

  \exercise{1} Compute $\Srel\rho\sigma$ for $\rho=\ket+\!\bra+$ and $\sigma=\Id/2$, for $\rho=\Id/2$ and $\sigma=\ket0\!\bra0$, and for $\rho=\diag(\tfrac12,\tfrac12)$, $\sigma=\diag(\tfrac14,\tfrac34)$.

  \exercise{2} Repeat Example~\ref{ex:qre}(d) with $\sigma=\tfrac12\ket+\!\bra++\tfrac12\ket-\!\bra-=\Id/2$ and with $\sigma=\tfrac{9}{10}\ket+\!\bra++\tfrac1{10}\ket-\!\bra-$, and compare.

  \exercise{2} Show that $\Srel\rho\sigma$ is not symmetric by computing it both ways for $\rho=\diag(\tfrac34,\tfrac14)$ and $\sigma=\diag(\tfrac12,\tfrac12)$.

  \exercise{2} Verify Corollary~\ref{cor:vn-consequences}(c) (concavity) for $\rho_1=\ket0\!\bra0$, $\rho_2=\ket+\!\bra+$ and $p=(\tfrac12,\tfrac12)$ by computing both sides explicitly.

  \exercise{2} Use the formula $\log(\rho\otimes\sigma)=\log\rho\otimes\Id+\Id\otimes\log\sigma$ to prove that $\Srel{\rho_1\otimes\rho_2}{\sigma_1\otimes\sigma_2}=\Srel{\rho_1}{\sigma_1}+\Srel{\rho_2}{\sigma_2}$.

  \exercise{3} Show that $\Srel{U\rho U^\dagger}{U\sigma U^\dagger}=\Srel\rho\sigma$ for every unitary $U$ (unitary invariance), and deduce that distinguishability is the same in any basis.

  \exercise{2} Verify the subadditivity inequality for $\bell00$, for $\ket+\ket0$, and for the state $\tfrac1{\sqrt3}(\ket{00}+\ket{01}+\ket{10})$ (use Example~\ref{ex:ent-entropy}).

  \exercise{3} Show that $I(A\!:\!B)\le2\min\{\Ss(\rho^A),\Ss(\rho^B)\}$ for a pure state of $AB$, and that equality holds for all pure states. For a general mixed state the bound is $I\le2\min\{\log d_A,\log d_B\}$.

  \exercise{3} For $\ket\psi=\cos\frac\theta2\ket0+\sin\frac\theta2\ket1$, compute $\Ss(\mathcal D(\ket\psi\!\bra\psi))$ for the dephasing in the computational basis, as a function of $\theta$, and check that it is $\ge0=\Ss(\ket\psi\!\bra\psi)$.

  \exercise{3} Let $\rho_0=\ket0\!\bra0$ and $\rho_1=\ket+\!\bra+$ be sent with probability $\tfrac12$ each. Compute the \emph{Holevo quantity} $\chi=\Ss(\bar\rho)-\tfrac12\Ss(\rho_0)-\tfrac12\Ss(\rho_1)$ and compare it with $1$ bit. (It bounds the information that Bob can extract about the bit sent.)
\end{exerciseset}
